% !TeX root = ../informe.tex
% ===========================================================================
%  Datos de la portada y metadatos del PDF.
%  Este es el único archivo que hay que editar con los datos del trabajo.
% ===========================================================================

\newcommand{\tituloInforme}{Título del informe}
\newcommand{\subtituloInforme}{Subtítulo o alcance del trabajo}
\newcommand{\nombreCurso}{Nombre del curso}
\newcommand{\nombrePrograma}{Nombre del programa}
\newcommand{\nombreUniversidad}{Nombre de la institución}
\newcommand{\tipoTrabajo}{Tipo de trabajo}

% --- Integrantes del equipo ------------------------------------------------
% Un par nombre/correo por integrante. Dejar vacíos los que sobren: la
% portada omite automáticamente los que estén en blanco.
\newcommand{\autorUno}{Nombre del primer autor}
\newcommand{\correoUno}{correo@ejemplo.org}

\newcommand{\autorDos}{}
\newcommand{\correoDos}{}

\newcommand{\autorTres}{}
\newcommand{\correoTres}{}

\newcommand{\autorCuatro}{}
\newcommand{\correoCuatro}{}

\newcommand{\nombreDocente}{Nombre del docente}
\newcommand{\ciudadFecha}{Ciudad, \today}

% --- Metadatos del PDF -----------------------------------------------------
% pdfauthor con varios autores: sepáralos con comas dentro de la clave.
\hypersetup{
    pdftitle  = {\tituloInforme},
    pdfauthor = {\autorUno},
    pdfsubject= {\nombreCurso\ --- \nombrePrograma},
    pdfkeywords = {palabra clave, palabra clave, palabra clave}
}
